\section{Sistematika Penulisan}

Sistematika penulisan laporan Proyek Akhir ini disusun secara sistematis agar pembaca dapat memahami alur pembahasan dengan jelas dan runtut. Penulisan dibagi ke dalam beberapa bab, mulai dari pendahuluan hingga lampiran, dengan rincian sebagai berikut:

\begin{longtable}{@{}llp{11cm}@{}}
    BAB I          & : & PENDAHULUAN                                                                                                                                                                                                                                                                                                                                                                                                                                                                                                                                                                                                                                                                                                                                                                                                            \\
                   &   & Bab ini memberikan gambaran umum mengenai latar belakang permasalahan yang melatarbelakangi dilakukannya proyek akhir, perumusan masalah yang akan dikaji, tujuan yang hendak dicapai, serta manfaat yang diharapkan baik secara teoritis maupun praktis. Selain itu, bab ini juga memuat batasan masalah yang menjadi ruang lingkup kajian agar fokus pembahasan lebih terarah. Di akhir bab, dijelaskan pula sistematika penulisan yang memberikan panduan mengenai struktur dan isi masing-masing bab dalam laporan ini.                                                                                                                                                                                                                                                                                            \\[0.8em]

    BAB II         & : & STUDI PUSTAKA                                                                                                                                                                                                                                                                                                                                                                                                                                                                                                                                                                                                                                                                                                                                                                                                          \\
                   &   & Bab ini membahas berbagai teori, konsep, dan referensi dari penelitian atau sumber ilmiah terdahulu yang relevan dengan topik proyek akhir. Kajian pustaka ini mencakup pembahasan mengenai sistem e-ticketing, mekanisme reservasi daring, pemanfaatan teknologi QRCode, serta ulasan terhadap perangkat teknologi yang digunakan dalam pengembangan sistem, baik dari sisi perangkat lunak maupun perangkat keras. Studi pustaka ini menjadi dasar pijakan dalam merumuskan solusi dan pendekatan teknis yang digunakan dalam proyek akhir ini.                                                                                                                                                                                                                                                                      \\[0.8em]

    BAB III        & : & METODE PENELITIAN                                                                                                                                                                                                                                                                                                                                                                                                                                                                                                                                                                                                                                                                                                                                                                                                      \\
                   &   & Bab ini menjelaskan pendekatan dan metode yang digunakan selama pelaksanaan proyek akhir. Penjelasan mencakup tahapan-tahapan dalam pengembangan sistem, mulai dari perencanaan, perancangan, implementasi, hingga pengujian. Penggunaan metode Agile dengan pendekatan Kanban dijelaskan sebagai strategi manajemen proyek yang fleksibel dan adaptif. Selain itu, dibahas pula integrasi sistem secara teknis, seperti penerapan one-time token untuk penggabungan otentikasi OAuth2 milik Google dengan sistem login konvensional menggunakan email dan sandi, integrasi API dan webhook untuk sistem pembayaran otomatis, serta rancangan struktur data linked list berbasis timestamp sebagai mekanisme rotasi kunci pada sistem QRCode yang terintegrasi dengan perangkat IoT secara offline dan fault-tolerant. \\[0.8em]

    BAB IV         & : & HASIL DAN PEMBAHASAN                                                                                                                                                                                                                                                                                                                                                                                                                                                                                                                                                                                                                                                                                                                                                                                                   \\
                   &   & Bab ini menyajikan hasil dari proses implementasi sistem yang telah dirancang pada bab sebelumnya. Uraian dilakukan secara sistematis mulai dari hasil pengembangan antarmuka pengguna, proses backend, serta integrasi sistem secara keseluruhan. Hasil pengujian terhadap sistem turut disertakan dalam bab ini, meliputi User Acceptance Testing (UAT) untuk menilai kesesuaian sistem terhadap kebutuhan pengguna, Black Box Testing untuk menguji fungsionalitas sistem, serta Load Testing untuk mengevaluasi performa sistem dalam kondisi beban tinggi. Pembahasan disusun secara analitis untuk mengkaji sejauh mana sistem memenuhi tujuan dan spesifikasi yang telah dirumuskan sebelumnya.                                                                                                                 \\[0.8em]

    BAB V          & : & PENUTUP                                                                                                                                                                                                                                                                                                                                                                                                                                                                                                                                                                                                                                                                                                                                                                                                                \\
                   &   & Bab ini merupakan bagian akhir dari laporan proyek akhir yang berisi kesimpulan dari seluruh proses pelaksanaan proyek, mulai dari identifikasi masalah hingga hasil pengujian sistem. Kesimpulan disusun berdasarkan temuan yang diperoleh selama pelaksanaan proyek. Selain itu, bab ini juga memuat saran-saran yang dapat digunakan sebagai acuan bagi pengembangan sistem lanjutan atau penelitian selanjutnya yang relevan dengan topik ini, dengan harapan dapat menyempurnakan dan memperluas cakupan sistem yang telah dikembangkan.                                                                                                                                                                                                                                                                          \\[0.8em]

    DAFTAR PUSTAKA & : & DAFTAR PUSTAKA                                                                                                                                                                                                                                                                                                                                                                                                                                                                                                                                                                                                                                                                                                                                                                                                         \\
                   &   & Bagian ini memuat seluruh referensi dan sumber literatur yang digunakan sebagai landasan teori dan acuan teknis dalam penyusunan laporan proyek akhir. Daftar pustaka dapat berupa buku, jurnal ilmiah, artikel daring, dokumentasi perangkat lunak, maupun sumber lainnya yang kredibel dan relevan.                                                                                                                                                                                                                                                                                                                                                                                                                                                                                                                  \\[0.8em]

    LAMPIRAN       & : & LAMPIRAN                                                                                                                                                                                                                                                                                                                                                                                                                                                                                                                                                                                                                                                                                                                                                                                                               \\
                   &   & Bagian lampiran menyajikan dokumen-dokumen pendukung yang tidak dicantumkan dalam bab utama, namun memiliki relevansi penting terhadap keseluruhan isi laporan. Lampiran dapat mencakup tangkapan layar sistem, potongan kode program penting, dokumentasi pengujian, konfigurasi sistem, diagram tambahan, serta dokumen lainnya yang mendukung validitas dan keterlacakan hasil proyek akhir.                                                                                                                                                                                                                                                                                                                                                                                                                        \\
\end{longtable}
